% =====================================================================================
%  Chapter 6 — Two-Robot Experimental Results
%  Thesis: Comparison of Control Strategies for Interaction-Force Aware Multirotor Teams
%
%  DRAFT SKELETON, 27 September 2026.
%  REVISED SKELETON 2026-10-14 (thesis-revision run, Phase 4; plan rows P6-01..P6-06). Structure
%  only: no number, figure or ranking from any flight, lab session, simulation or replay appears
%  here. Every results slot is a \todo that names the table or figure, the statistic, and what a
%  negative answer would look like.
%  Sources: Chapter~\ref{ch:setup} (platform, metrics, statistics), Chapter~\ref{ch:control}.
%
%  HONESTY NOTE. This file is structure and framing only. No comparison numbers, figures,
%  or rankings appear here. Every slot for measured outcomes is marked with \todo so the compiled
%  PDF cannot be read as reporting results that do not yet exist.
%
%  Citation baseline: 2026-08-27 (see docs/thesis/CITATION_BASELINE.md).
% =====================================================================================

\chapter{Two-Robot Experimental Results}
\label{ch:results2}

% CHANGED (P6-02): internal campaign label removed.
This chapter reports the systematic comparison on the two-vehicle team. Chapter~\ref{ch:setup}
fixed the platform, scenarios, metrics and protocol; Chapter~\ref{ch:control}
defined Strategies~0--3 compared against the baseline. The sections below follow the analysis
from aggregate tracking metrics through control effort, the predictive residual path, and the
attribution of effects to a loop, to a summary of when each strategy is preferable. Research
questions RQ1, RQ2 and RQ4 (Section~\ref{sec:intro-rq}) are answered primarily here; RQ3 is
deferred to the multi-robot results chapter.

All metrics are those defined in Section~\ref{sec:setup-metrics}, computed per flight and summarised
per (scenario, strategy) cell as specified in Section~\ref{sec:setup-stats}.

% -------------------------------------------------------------------------------------
\section{Experimental campaign overview}
\label{sec:res2-overview}

% DECISION D-05 (H9): the three paragraphs of this section name vehicles by nickname, the camera
%   count, internal campaign labels (C.1/C.4) and the current campaign status ("Indicative single
%   flights and partial C.1 collection exist"). Options: (i) replace by a neutral description to be
%   completed after the campaign (recommended: platform as in Chapter 5, scenarios flown, n per
%   cell, deviations from the protocol); (ii) keep and update at the end. Text left unchanged.
The comparison flights use the brushless pair \textbf{cf5} (study vehicle, bottom role in vertical
stacks) and \textbf{cf\_second} (top), with external pose from a four-camera OptiTrack system and
onboard logging at \SI{500}{\hertz} on micro-SD as the source of record (Section~\ref{sec:setup-platform}).
Radio telemetry is for monitoring only; thesis metrics are computed from merged uSD logs after role
alignment, using the same pipeline documented in Chapter~\ref{ch:setup}.

Scenarios are drawn from the frozen formation library (Chapter~\ref{ch:setup}, Section~\ref{sec:setup-scenarios}):
the C.1/C.4 matrix prioritises vertical separation and downwash-dominated geometries (e.g.\ A1,
A3, A4, A7) together with null coplanar controls (C1--C3) that bound the measurement floor. Each
(scenario, strategy) cell is intended to be repeated $n$ times under interleaved runtime controller
selection so session effects do not align with strategy choice (Table~\ref{tab:setup-threats}).

The compared set follows Chapter~\ref{ch:control}: Strategy~0 (uncompensated geometric baseline),
Strategy~1 (pure INDI), Strategy~2 (geometric position loop with onboard Neural-Swarm2-style
residual prediction), and Strategy~3 (feedback linearisation with the same predictor), each flown
with gains frozen after the hardware gate. Indicative single flights and partial C.1 collection
exist; the tables and figures in the following sections are reserved for the completed C.4 repeats.

% NEW (P6-02): campaign table slot.
\todo[color=red!25]{PLACEHOLDER: Table --- scenarios flown $\times$ strategies, with $n$ per cell,
flights excluded and the reason for each exclusion, and every deviation from the protocol of
Section~\ref{sec:setup-protocol}.}

% -------------------------------------------------------------------------------------
\section{Tracking and attitude}
\label{sec:res2-tracking}

This section answers the tracking half of RQ1: whether compensating strategies reduce position and
separation error relative to Strategy~0 under identical trajectories, and how reactive versus
predictive families rank. Attitude excursions matter for safety and for interpreting residual
magnitude; they are reported alongside position error rather than as a separate optimisation target.

% CHANGED (P6-02): metrics as defined in Chapter 5; statistic; negative answer.
\todo[color=red!25]{PLACEHOLDER: Table~\ref{tab:res2-tracking}. Columns: scenario, strategy,
position RMSE per axis and total over the window $W$, maximum deviation, mean and standard deviation
of the separation error $e_{\mathrm{sep}}$, RMSE and peak of $\lVert e_R\rVert$, remaining
$\operatorname{RMS}_W(a_{\mathrm{res}})$, $n$. Rows: all (scenario, strategy) cells flown. Statistic:
the per-cell summary and the paired within-session difference to Strategy~0 fixed in
Section~\ref{sec:setup-stats}, judged against the measurement floor of the null scenarios by the
rule fixed there.}

\begin{table}[t]
  \centering
  \caption{Tracking error summary (two-robot comparison campaign).}
  \label{tab:res2-tracking}
  \small
  \begin{tabular}{@{}llrrrrr@{}}
    \toprule
    Scenario & Strategy & RMSE $x$ & RMSE $y$ & RMSE $z$ & Peak $\|\mathbf{e}\|$ & $n$ \\
    \midrule
    \multicolumn{7}{@{}l@{}}{\textit{--- populated from the per-flight metrics ---}} \\
    \makecell{Placeholder} & Strategy~0 & --- & --- & --- & --- & --- \\
    \bottomrule
  \end{tabular}
\end{table}

\todo[color=red!25]{PLACEHOLDER: Figure (grouped bars) --- strategies $\times$ metrics, one panel per
metric family (position RMSE, separation error, attitude error), error bars as fixed in
Section~\ref{sec:setup-stats}.}

\todo[color=red!25]{PLACEHOLDER: Figure (time series) --- commanded and achieved separation for one
representative transient scenario (A8 or A4), Strategies~0--2 on the same axes, crossing time from
the commanded geometry marked.}

% CHANGED (P6-06, H5): only Strategy 1 is reported as INDI.
A short subsection will contrast geometric tracking (Strategy~0) with the INDI implementation that
constitutes Strategy~1 (Section~\ref{sec:ctrl-s1}). No other INDI controller is reported in this
comparison.

% NEW (P6-02): what a negative answer looks like.
% CHANGED (closing pass, diff audit): negative answers refer to the decision rule fixed in Sec. setup-stats
%   instead of stating one.
\paragraph{Negative answer.} If no compensating strategy differs from Strategy~0 beyond the
measurement floor on the separation error (by the rule of Section~\ref{sec:setup-stats}), the answer to the tracking half of RQ1 is that, at
the separations flown, compensation does not measurably improve tracking, and the section reports
that rather than a ranking.

% -------------------------------------------------------------------------------------
\section{Control effort and safety}
\label{sec:res2-effort}

Better tracking is not automatically preferable if it requires sustained saturation or leaves less
margin on separation. This section supports RQ4 by pairing performance with effort proxies and
closest-approach statistics.

% CHANGED (P6-02): internal variable names removed; metrics as defined in Chapter 5.
\todo[color=red!25]{PLACEHOLDER: Table~\ref{tab:res2-effort} --- RMS commanded torque, RMS change of
the motor commands, fraction of samples with a motor command at its limit, minimum achieved
horizontal and vertical separation; per-cell summary and $n$ as in Section~\ref{sec:setup-stats}.}

\begin{table}[t]
  \centering
  \caption{Control effort and safety margins (two-robot comparison campaign).}
  \label{tab:res2-effort}
  \small
  \begin{tabular}{@{}llrrrr@{}}
    \toprule
    Scenario & Strategy & Effort proxy & Sat.\ fraction & Min sep.\ (\si{\metre}) & $n$ \\
    \midrule
    \makecell{Placeholder} & Strategy~0 & --- & --- & --- & --- \\
    \bottomrule
  \end{tabular}
\end{table}

The prose here will state explicitly when a strategy trades higher effort for lower tracking error,
and when reduced separation margin disqualifies a configuration despite good RMSE.

% -------------------------------------------------------------------------------------
\section{Predictive strategy (Strategy~2)}
\label{sec:res2-nn}

% CHANGED (P6-05): the two kinds of flight defined in Chapter 5, Sec. setup-campaign.
Strategy~2 is evaluated in two layers. On \emph{prediction-off} flights, flown under Strategy~0 with
the trained weights loaded, the prediction is recorded but does not drive the loop; these flights
measure model quality as the prediction error $\epsilon_{\mathrm{pred}}$ of
Section~\ref{sec:setup-metrics}, independent of loop gain. On \emph{prediction-on} flights the
prediction drives the controller, and these are compared with Strategy~0 on matched scenarios. RQ1
and RQ2 both depend on the second; the first explains it.

\todo[color=red!25]{PLACEHOLDER: Figure --- predicted against measured $a_{\mathrm{res},z}$ on
prediction-off flights (scatter, with bias and RMSE); second panel: error against relative geometry
(lateral distance, $\Delta z$, whether the neighbour is inside the input gate). Training scenarios
held out from each panel stated in the caption.}

\todo[color=red!25]{PLACEHOLDER: Table --- Strategy~2 (prediction-on) against Strategy~0 on matched
scenarios, same columns as Table~\ref{tab:res2-tracking}; caption states the training scenarios and
the weight set used.}

\paragraph{Negative answer.} A model whose prediction error on prediction-off flights is not smaller
than the measured residual itself cannot be expected to help, and a prediction-on improvement
within the measurement floor (Section~\ref{sec:setup-stats}) is reported as no improvement.

% NEW (P6-04): the checks promised in Chapter 3 (Sec. model-sign) and Chapter 4 (Sec. ctrl-s1).
\section{Sign and loop attribution}
\label{sec:res2-attribution}

Two checks connect the measured differences to the mechanisms of Chapters~\ref{ch:model}
and~\ref{ch:control}. First, the sign of the compensation: Chapter~\ref{ch:model} predicts that with
the sign correct the formation-induced displacement in a quasi-static stack is smaller than under
Strategy~0, and with the sign inverted approximately twice as large. Second, the loop that carries
an effect: the position-only and attitude-only INDI configurations of Section~\ref{sec:ctrl-s1} are
diagnostic, not strategies, and are used only to attribute a Strategy~1 effect to the position or
the attitude increment.
% DECISION D-01: which sign Strategy 1 flies determines which outcome of the first check counts as
%   "correct"; this section cannot be written until D-01 is decided.

\todo[color=red!25]{PLACEHOLDER: Table --- quasi-static stack (A1), separation error of Strategy~1
against Strategy~0, with the expected direction stated before the data; diagnostic configurations
reported separately and labelled as diagnostic.}

% REMOVED (P6-03): the former section "Hybrid (NA-INDI), if in scope", which described a hybrid as
%   "Strategy 4 in the wider project numbering". No strategy above 3 belongs in the main chapters
%   (Chapter 4 scope); the hybrid is discussed in Chapter 2 and is not compared. Original text:
%   "The primary compared set in Chapter~\ref{ch:control} is Strategies~0--3. A NA-INDI / LINDI
%   hybrid (Strategy~4 in the wider project numbering) is \emph{not} in that set unless explicitly
%   restored to the C.4 matrix. This section remains as a structured slot: if hybrid flights are
%   flown, they use the same metric tables as Sections~\ref{sec:res2-tracking}--\ref{sec:res2-effort}
%   for fair comparison." The label is kept below so that existing references resolve.
\label{sec:res2-hybrid}

% -------------------------------------------------------------------------------------
\section{Summary --- when which strategy}
\label{sec:res2-summary}

This section closes the two-robot story for RQ1 and RQ2: a scenario-class breakdown (vertical stack,
lateral crossing, near-ground, fast transient) with both qualitative ranking and the quantitative
backing from Tables~\ref{tab:res2-tracking}--\ref{tab:res2-effort}.

\todo[color=red!25]{PLACEHOLDER: Summary table --- rows: strategies; columns: scenario classes;
cells: best tracking / effort trade-off / data and compute cost (RQ4 pointers). No ranking is
entered without the statistical support of Section~\ref{sec:setup-stats}.}

% DECISION D-05 (H5): the Strategy 3 row describes an unimplemented strategy as compared; keep only
%   if Strategy 3 is implemented, otherwise mark "not implemented". Row left unchanged.
\begin{table}[t]
  \centering
  \caption{Strategy suitability by scenario class (to be filled after the comparison campaign).}
  \label{tab:res2-when-which}
  \small
  \begin{tabular}{@{}lcccc@{}}
    \toprule
    Strategy & Vertical stack & Lateral traverse & Near-ground & Fast transient \\
    \midrule
    Strategy~0 (baseline) & --- & --- & --- & --- \\
    Strategy~1 (INDI) & --- & --- & --- & --- \\
    Strategy~2 (pred.\ + geo) & --- & --- & --- & --- \\
    Strategy~3 (pred.\ + FBL) & --- & --- & --- & --- \\
    \bottomrule
  \end{tabular}
\end{table}

Train/test splits by scenario class (Section~\ref{sec:setup-validity}) feed RQ2 directly: a crossover
in favour of reactive or predictive compensation should appear here if it exists in the data.
% CHANGED (closing pass, diff audit): marker added; rule referenced, not stated.
\paragraph{Negative answer.} If no strategy is preferable in one class and not in another beyond the
measurement floor (Section~\ref{sec:setup-stats}), RQ2 is answered as ``no regime dependence detected at the separations and
speeds flown'', with the excitation actually achieved stated next to it.
